\documentclass[12pt,a4paper]{article}

% --- Codificação e língua ---
\usepackage[utf8]{inputenc}
\usepackage[T1]{fontenc}
\usepackage[brazil]{babel}

% --- Layout ---
\usepackage[top=2.5cm, bottom=2.5cm, left=3cm, right=2.5cm]{geometry}
\usepackage{indentfirst}
\setlength{\parskip}{6pt}

% --- Tipografia ---
\usepackage{lmodern}
\usepackage{microtype}

% --- Matemática ---
\usepackage{amsmath}
\usepackage{amssymb}

% --- Código-fonte ---
\usepackage{listings}
\usepackage{xcolor}

\definecolor{codebg}{RGB}{245,245,245}
\definecolor{codecomment}{RGB}{100,100,100}
\definecolor{codestring}{RGB}{163,21,21}
\definecolor{codekeyword}{RGB}{0,0,200}

\lstset{
  backgroundcolor=\color{codebg},
  basicstyle=\ttfamily\small,
  breaklines=true,
  breakatwhitespace=false,
  commentstyle=\color{codecomment}\itshape,
  keywordstyle=\color{codekeyword}\bfseries,
  stringstyle=\color{codestring},
  showstringspaces=false,
  frame=single,
  framesep=4pt,
  rulecolor=\color{gray!40},
  numbers=none,
  captionpos=b,
  tabsize=4,
}

% Estilo especial para blocos de shell/saída (sem highlight)
\lstdefinestyle{shell}{
  language={},
  basicstyle=\ttfamily\small,
  backgroundcolor=\color{codebg},
  frame=single,
  framesep=4pt,
  rulecolor=\color{gray!40},
  breaklines=true,
}

% Estilo para C
\lstdefinestyle{ccode}{
  language=C,
  basicstyle=\ttfamily\small,
  backgroundcolor=\color{codebg},
  frame=single,
  framesep=4pt,
  rulecolor=\color{gray!40},
  keywordstyle=\color{codekeyword}\bfseries,
  commentstyle=\color{codecomment}\itshape,
  stringstyle=\color{codestring},
  breaklines=true,
}

% --- Tabelas ---
\usepackage{booktabs}
\usepackage{array}

% --- Imagens ---
\usepackage{graphicx}

% --- Hiperlinks ---
\usepackage[colorlinks=true, linkcolor=black, urlcolor=blue, citecolor=black]{hyperref}

% --- Títulos de seção ---
\usepackage{titlesec}
\titleformat{\section}{\large\bfseries}{}{0em}{}[\titlerule]
\titleformat{\subsection}{\normalsize\bfseries}{}{0em}{}

% --- Espaçamento de linha ---
\usepackage{setspace}
\onehalfspacing

% =========================================================
\begin{document}

% ----- CAPA -----
\begin{titlepage}
  \centering
  \vspace*{3cm}
  {\LARGE\bfseries Relatório Prático\\[0.4em]
   Programação Paralela com MPI\par}
  \vspace{2cm}
  {\large
    \textbf{Nome:} Vivian Souza \\[0.5em]
    \textbf{DRE:} 123205793 \\[0.5em]
    \textbf{Data:} 19 de Junho de 2026
  }
  \vfill
\end{titlepage}

\tableofcontents
\newpage

% =========================================================
\section{Exercício 1 — Variantes do Hello World}

\subsection{1(a) Introspecção em Nível de SO}

Cada processo MPI foi instrumentado com as chamadas \texttt{getpid()} (de \texttt{<unistd.h>}) e \texttt{sched\_getcpu()} (de \texttt{<sched.h>}) para reportar, além do seu rank, o PID real no Sistema Operacional e o núcleo físico em que estava sendo executado no momento da impressão. O programa foi compilado com a diretiva \texttt{\#define \_GNU\_SOURCE} para habilitar a função \texttt{sched\_getcpu()}.

\textbf{Execução:}
\begin{lstlisting}[style=shell]
mpicc -O2 -Wall -o hello_os hello_os.c
mpiexec -n 4 ./hello_os
mpiexec -n 8 --oversubscribe ./hello_os
\end{lstlisting}

\textbf{Saídas observadas (múltiplas execuções com \texttt{-n 4}):}

\begin{lstlisting}[style=shell]
Run 1:
Hello from rank 0/4 -- PID = 29237, CPU = 3
Hello from rank 1/4 -- PID = 29238, CPU = 7
Hello from rank 2/4 -- PID = 29239, CPU = 1
Hello from rank 3/4 -- PID = 29242, CPU = 5

Run 2:
Hello from rank 0/4 -- PID = 29255, CPU = 0
Hello from rank 1/4 -- PID = 29256, CPU = 2
Hello from rank 2/4 -- PID = 29257, CPU = 1
Hello from rank 3/4 -- PID = 29259, CPU = 3

Run 3:
Hello from rank 0/4 -- PID = 29285, CPU = 0
Hello from rank 1/4 -- PID = 29286, CPU = 2
Hello from rank 2/4 -- PID = 29287, CPU = 3
Hello from rank 3/4 -- PID = 29289, CPU = 4
\end{lstlisting}

Também foram feitas execuções com \texttt{-n 8} e \texttt{--oversubscribe}, pois o computador não dispunha de núcleos físicos suficientes para executar sem essa flag. Saídas observadas:

\begin{lstlisting}[style=shell]
Run 1:
Hello from rank 0/8 -- PID = 33737, CPU = 4
Hello from rank 1/8 -- PID = 33738, CPU = 3
Hello from rank 2/8 -- PID = 33739, CPU = 6
Hello from rank 3/8 -- PID = 33740, CPU = 2
Hello from rank 4/8 -- PID = 33741, CPU = 1
Hello from rank 5/8 -- PID = 33745, CPU = 5
Hello from rank 6/8 -- PID = 33747, CPU = 3
Hello from rank 7/8 -- PID = 33749, CPU = 0

Run 2:
Hello from rank 0/8 -- PID = 33821, CPU = 4
Hello from rank 1/8 -- PID = 33822, CPU = 7
Hello from rank 2/8 -- PID = 33823, CPU = 0
Hello from rank 3/8 -- PID = 33824, CPU = 6
Hello from rank 4/8 -- PID = 33826, CPU = 3
Hello from rank 5/8 -- PID = 33829, CPU = 1
Hello from rank 6/8 -- PID = 33832, CPU = 5
Hello from rank 7/8 -- PID = 33835, CPU = 2

Run 3:
Hello from rank 0/8 -- PID = 33869, CPU = 4
Hello from rank 1/8 -- PID = 33870, CPU = 6
Hello from rank 2/8 -- PID = 33871, CPU = 1
Hello from rank 3/8 -- PID = 33872, CPU = 7
Hello from rank 4/8 -- PID = 33875, CPU = 0
Hello from rank 5/8 -- PID = 33877, CPU = 2
Hello from rank 6/8 -- PID = 33880, CPU = 3
Hello from rank 7/8 -- PID = 33882, CPU = 5
\end{lstlisting}

\textbf{Observações:}

As execuções sucessivas confirmam na prática dois princípios fundamentais da arquitetura SPMD do MPI:

\begin{enumerate}
  \item \textbf{Todos os PIDs são diferentes entre execuções:} o MPI não opera com threads compartilhadas, mas com processos reais e isolados no Sistema Operacional. Cada rank possui seu próprio espaço de endereçamento de memória e um PID exclusivo atribuído pelo kernel. A cada nova invocação de \texttt{mpiexec}, um conjunto inteiramente novo de PIDs é gerado, independentemente do valor de $p$.

  \item \textbf{A coluna de CPU varia entre execuções:} o ambiente MPI delega ao escalonador dinâmico do kernel do Linux a decisão de qual núcleo físico processará cada tarefa. Como o escalonamento depende da carga corrente da máquina e não é determinístico, os processos ``saltam'' de núcleo a cada rodada. Um ponto interessante observado com \texttt{--oversubscribe} é que diferentes ranks/threads podem ser alocados na mesma CPU.
\end{enumerate}

% ---------------------------------------------------------
\subsection{1(b) Variante em Anel}

Na topologia em anel, o processo de rank $k$ envia sua saudação ao processo $(k+1) \bmod p$ e recebe do processo $(k-1+p) \bmod p$. A expressão \texttt{(k-1+p) \% p} é necessária pois, em C, o operador \texttt{\%} pode retornar valores negativos quando o dividendo é negativo: para $k=0$, \texttt{(0-1) \% p} resultaria em $-1$ em vez do esperado $p-1$. Somar $p$ antes da operação de módulo garante sempre um argumento não negativo.

Uma abordagem ingênua consistiria em todos os processos chamarem \texttt{MPI\_Recv} seguido de \texttt{MPI\_Send}. Como \texttt{MPI\_Recv} é estritamente bloqueante, todos os processos ficariam indefinidamente aguardando receber uma mensagem, sem que nenhum deles chegasse à linha de \texttt{MPI\_Send}, criando um ciclo de espera infinita (deadlock).

Uma das alternativas seria fazer, por exemplo, o rank 0 executar \texttt{MPI\_Send} antes de \texttt{MPI\_Recv}, enquanto os demais fariam \texttt{MPI\_Recv} antes de \texttt{MPI\_Send}, quebrando a simetria do ciclo. Apesar de funcionar em alguns casos, essa abordagem depende do comportamento interno da implementação do MPI, já que o envio pode bloquear caso a mensagem não caiba nos buffers disponíveis. Além disso, o código deixa de ser simétrico, tornando a implementação mais difícil de manter.

Por esse motivo, a solução adotada foi o uso de \texttt{MPI\_Sendrecv}, que foi desenvolvido justamente para esse tipo de comunicação. A função realiza o envio e o recebimento de forma segura, evitando deadlocks independentemente do tamanho da mensagem ou da quantidade de processos. Outra vantagem é que todos os ranks executam exatamente o mesmo código, preservando o modelo SPMD e deixando a implementação mais simples, legível e portável.

\textbf{Execução:}
\begin{lstlisting}[style=shell]
mpicc -O2 -Wall -o hello_ring hello_ring.c
mpiexec -n 4 ./hello_ring
mpiexec -n 8 --oversubscribe ./hello_ring
\end{lstlisting}

\textbf{Saída:}
\begin{lstlisting}[style=shell]
Rank 0 recebeu a mensagem: Greeting from rank 3 of 4! (do Rank 3)
Rank 0 coletou do Rank 1 a mensagem: Greeting from rank 0 of 4!
Rank 0 coletou do Rank 2 a mensagem: Greeting from rank 1 of 4!
Rank 0 coletou do Rank 3 a mensagem: Greeting from rank 2 of 4!

Rank 0 recebeu a mensagem: Greeting from rank 7 of 8! (do Rank 7)
Rank 0 coletou do Rank 1 a mensagem: Greeting from rank 0 of 8!
Rank 0 coletou do Rank 2 a mensagem: Greeting from rank 1 of 8!
Rank 0 coletou do Rank 3 a mensagem: Greeting from rank 2 of 8!
Rank 0 coletou do Rank 4 a mensagem: Greeting from rank 3 of 8!
Rank 0 coletou do Rank 5 a mensagem: Greeting from rank 4 of 8!
Rank 0 coletou do Rank 6 a mensagem: Greeting from rank 5 of 8!
Rank 0 coletou do Rank 7 a mensagem: Greeting from rank 6 of 8!
\end{lstlisting}

A saída valida o tráfego correto: o rank 0 recebeu de 3 (seu vizinho à esquerda no anel), e ao longo das $p-1$ iterações colecionou todas as mensagens originadas pelos demais ranks, em ordem crescente de rank de origem, conforme o percurso circular.

% =========================================================
\section{Exercício 2 — Regra do Trapézio Generalizada}

\subsection{2(a) Balanceamento de Carga para $n$ Arbitrário}

O código base \texttt{mpi\_trap\_input.c} da Aula 5 foi estendido para lidar com valores de $n$ que não são divisíveis pelo número de processos $p$. A estratégia adotada distribui o resto da divisão entre os primeiros ranks: os $n \bmod p$ primeiros processos recebem $\lceil n/p \rceil$ trapézios, e os demais recebem $\lfloor n/p \rfloor$. O bloco de cálculo dos parâmetros locais ficou:

\begin{lstlisting}[style=ccode]
int rem      = n % comm_sz;
int local_n  = n / comm_sz + (my_rank < rem ? 1 : 0);
int offset   = my_rank * (n / comm_sz) + (my_rank < rem ? my_rank : rem);
double local_a = a + offset * h;
double local_b = local_a + local_n * h;
\end{lstlisting}

\textbf{Verificação com \texttt{p=3}, \texttt{n=10}, $f(x)=\sin(x)$ em $[0,\pi]$ (valor exato~$= 2$):}

\begin{lstlisting}[style=shell]
echo "0 3.141592653589793 10" | mpiexec -n 3 ./mpi_trap_generalized
\end{lstlisting}

\textbf{Saída:}
\begin{lstlisting}[style=shell]
Enter a, b, n: Rank 0 executou 4 trapezios no intervalo [0.000000, 1.256637]
Rank 1 executou 3 trapezios no intervalo [1.256637, 2.199115]
Rank 2 executou 3 trapezios no intervalo [2.199115, 3.141593]
With n = 10 trapezoids, integral of f on [0, 3.14159] = 1.983523537509454e+00
\end{lstlisting}

A distribuição 4+3+3 está correta: $10 \bmod 3 = 1$, portanto apenas o rank 0 recebe um trapézio a mais. O resultado numérico $\approx 1{,}9835$ concorda com a aproximação esperada para $n=10$ pontos na regra do trapézio aplicada a $\sin(x)$.

% ---------------------------------------------------------
\subsection{2(b) Estudo de Escalabilidade Forte ($n=10^7$)}

O cronômetro foi posicionado imediatamente após \texttt{MPI\_Init}, de modo que $T_p$ inclui o overhead de \texttt{Get\_input}, o bloco de computação e a comunicação de coleta — constituindo uma referência justa de processo único conforme solicitado pelo enunciado. O \texttt{MPI\_Barrier} da versão anterior foi removido por ser um artefato desnecessário. O tempo paralelo reportado é o \emph{máximo} entre todos os processos, obtido via \texttt{MPI\_Reduce} com \texttt{MPI\_MAX}, refletindo o tempo de relógio real da execução.

\textbf{Execução:}
\begin{lstlisting}[style=shell]
echo "0 3.141592653589793 10000000" | mpiexec -n 1 ./mpi_trap_generalized
echo "0 3.141592653589793 10000000" | mpiexec -n 2 ./mpi_trap_generalized
echo "0 3.141592653589793 10000000" | mpiexec -n 4 ./mpi_trap_generalized
echo "0 3.141592653589793 10000000" | mpiexec -n 8 --oversubscribe ./mpi_trap_generalized
\end{lstlisting}

\textbf{Tabela de desempenho} ($T_s = 0{,}117264$~s):

\begin{table}[h!]
\centering
\begin{tabular}{@{}ccccc@{}}
\toprule
$p$ & $T_p$ (s) & $S_p = T_s/T_p$ & $E_p = S_p/p$ \\
\midrule
1 & 0{,}117264 & 1{,}000 & 1{,}000 \\
2 & 0{,}066322 & 1{,}768 & 0{,}884 \\
4 & 0{,}043899 & 2{,}671 & 0{,}668 \\
8 & 0{,}026557 & 4{,}416 & 0{,}552 \\
\bottomrule
\end{tabular}
\caption{Speedup e eficiência para $n=10^7$.}
\end{table}

\textbf{Discussão — a partir de qual $p$ a eficiência $E_p$ cai abaixo de 0,8?}

A eficiência cai abaixo de 0,8 a partir de $p=4$, atingindo $E_4 = 0{,}668$. A causa é multifatorial:

\begin{itemize}
  \item \textbf{Gargalo de comunicação linear:} a coleta dos resultados parciais é feita por um laço sequencial de chamadas \texttt{MPI\_Recv} na raiz (topologia em estrela). O custo dessa etapa cresce com $O(p)$. Com $p=4$ ou $p=8$, o tempo de cálculo da integral torna-se ínfimo, e o overhead de coordenação de mensagens passa a dominar o tempo total de execução.

  \item \textbf{Oversubscription (sobrecarga de CPU):} para $p=8$ foi necessária a flag \texttt{--oversubscribe}, indicando que a máquina física dispõe de menos núcleos reais do que os processos requisitados. O escalonador do SO é forçado a realizar trocas de contexto entre processos compartilhando o mesmo núcleo físico, o que degrada severamente o ganho paralelo e explica a eficiência de apenas 55{,}2\%.

  \item \textbf{Lei de Amdahl:} a fração sequencial irredutível do código (inicialização do ambiente MPI, leitura/escrita em tela) não é acelerada pelo paralelismo e impõe um limite superior teórico ao speedup, independentemente do número de processos.
\end{itemize}

% ---------------------------------------------------------
\subsection{2(c) Estudo de Convergência (Bônus)}

O programa paralelo foi executado com $p=4$ fixo, variando $n \in \{64, 256, 1024, 4096, 16384\}$. Para cada valor de $n$, o passo $h = \pi/n$ e o erro absoluto $|I_{\text{aprox}} - 2|$ foram registrados.

\textbf{Execução:}
\begin{lstlisting}[style=shell]
for n in 64 256 1024 4096 16384; do
    echo "0 3.141592653589793 $n" | mpiexec -n 4 ./mpi_trap_generalized
done
\end{lstlisting}

\textbf{Resultados:}

\begin{table}[h!]
\centering
\begin{tabular}{@{}cccl@{}}
\toprule
$n$ & $h = \pi/n$ & $I_{\text{aprox}}$ & $|I_{\text{aprox}} - 2|$ \\
\midrule
64    & 0{,}04909  & 1{,}999598388640037 & $4{,}016 \times 10^{-4}$ \\
256   & 0{,}01227  & 1{,}999974900235053 & $2{,}510 \times 10^{-5}$ \\
1024  & 0{,}00307  & 1{,}999998431268383 & $1{,}569 \times 10^{-6}$ \\
4096  & 0{,}000767 & 1{,}999999901954288 & $9{,}805 \times 10^{-8}$ \\
16384 & 0{,}000192 & 1{,}999999993872146 & $6{,}128 \times 10^{-9}$ \\
\bottomrule
\end{tabular}
\caption{Convergência da regra do trapézio para $f(x)=\sin(x)$ em $[0,\pi]$.}
\end{table}

A teoria do cálculo numérico prevê que o erro global da regra do trapézio decai como $O(h^2)$: quando $h$ é reduzido por um fator 4 (de $n=64$ para $n=256$), o erro deve cair por um fator $\approx 16$. Observa-se que ao passar de $n=64$ para $n=256$ o erro reduz de $4{,}016 \times 10^{-4}$ para $2{,}510 \times 10^{-5}$ — uma redução de fator $\approx 16$, confirmando a ordem quadrática.

Ao traçar $|I_{\text{aprox}} - 2|$ em função de $h$ em escala log-log, a regressão linear sobre os cinco pontos resulta em uma inclinação de \textbf{2,0000}, atestando inequivocamente a convergência quadrática teórica. A divisão de domínio e a coleta paralela via MPI não introduziram qualquer anomalia matemática, preservando integralmente as propriedades numéricas do algoritmo original.

\begin{figure}[h!]
  \centering
  \includegraphics[width=0.75\textwidth]{ex2_trap/grafico_convergencia.png}
  \caption{Gráfico log-log do erro absoluto em função de $h$. A inclinação da reta ajustada é $\approx 2{,}0$.}
\end{figure}

% =========================================================
\section{Exercício 3 — Soma Paralela (Ponto-a-Ponto)}

O programa foi estruturado em quatro etapas:

\begin{enumerate}
  \item O processo 0 gera um vetor de $N = 10^7$ doubles aleatórios (via \texttt{rand()}) e calcula a soma serial como gabarito.
  \item O processo 0 distribui blocos de $N/p$ elementos a cada processo trabalhador com um laço de \texttt{MPI\_Send}.
  \item Cada processo calcula sua soma local iterando sobre o bloco recebido.
  \item Os processos trabalhadores enviam suas somas parciais ao processo 0 via \texttt{MPI\_Send}; o processo 0 recebe com um laço de \texttt{MPI\_Recv} e acumula o total.
\end{enumerate}

\textbf{Execução:}
\begin{lstlisting}[style=shell]
mpiexec -n 1 ./mpi_psum
mpiexec -n 2 ./mpi_psum
mpiexec -n 4 ./mpi_psum
mpiexec -n 8 --oversubscribe ./mpi_psum
\end{lstlisting}

\textbf{Saída:}
\begin{lstlisting}[style=shell]
Soma Serial = 4.9998874970e+06
Soma Global = 4.9998874970e+06
Erro Relativo = 3.9488965689e-14

Soma Serial = 4.9984512878e+06
Soma Global = 4.9984512878e+06
Erro Relativo = 5.7573568088e-14

Soma Serial = 4.9996657464e+06
Soma Global = 4.9996657464e+06
Erro Relativo = 0.0000000000e+00

Soma Serial = 5.0002277570e+06
Soma Global = 5.0002277570e+06
Erro Relativo = 7.8227532894e-15
\end{lstlisting}

A soma global paralela concorda com o gabarito serial em todos os casos. O erro relativo residual (da ordem de $10^{-14}$ a $10^{-15}$) não representa um defeito no código: aritmética de ponto flutuante não é estritamente associativa, de modo que somar blocos em ordens diferentes altera os arredondamentos nos bits menos significativos. Com $p=1$ o erro é exatamente zero, pois os dois caminhos de soma são idênticos.

\textbf{Comparação com a Aula 6 (comunicação coletiva):} a implementação com primitivas ponto-a-ponto exigiu aproximadamente 18–20 linhas dedicadas apenas ao controle de comunicação: laço de \texttt{MPI\_Send} para distribuição, laço de \texttt{MPI\_Recv} para coleta, controle de índices e offsets de memória, e blocos \texttt{if/else} diferenciando mestre e trabalhadores. Com as rotinas coletivas da Aula 6, toda a lógica de distribuição colapsaria em uma única chamada \texttt{MPI\_Scatter}, e toda a coleta com redução colapsaria em uma única chamada \texttt{MPI\_Reduce} — substituindo $\sim$18 linhas por 2. Além da redução na contagem de linhas e na superfície exposta a bugs, as rotinas coletivas utilizam algoritmos internos em formato de árvore, substituindo a latência linear $O(p)$ da implementação manual por um custo assintótico $O(\log p)$.

% =========================================================
\section{Exercício 4 — Prática de Comunicação Coletiva}

\subsection{4(a) Hello World com \texttt{MPI\_Gather}}

Cada processo gerou sua saudação em um buffer local com \texttt{snprintf}. Em vez de um laço de \texttt{MPI\_Recv} no processo 0, foi utilizado \texttt{MPI\_Gather} para coletar todas as strings em um único vetor contíguo. O processo 0 alocou dinamicamente \texttt{comm\_sz * MAX\_STRING} bytes e, após o gather, percorreu o vetor com saltos de \texttt{MAX\_STRING} bytes para imprimir as mensagens em ordem de rank.

\textbf{Execução:}
\begin{lstlisting}[style=shell]
mpicc -O2 -Wall -o hello_gather hello_gather.c
mpiexec -n 4 ./hello_gather
\end{lstlisting}

\textbf{Saída:}
\begin{lstlisting}[style=shell]
Greetings from process 0 of 4!
Greetings from process 1 of 4!
Greetings from process 2 of 4!
Greetings from process 3 of 4!
\end{lstlisting}

A saída apresentou todas as saudações perfeitamente ordenadas por rank, confirmando que \texttt{MPI\_Gather} insere as contribuições de cada processo na posição correspondente ao seu rank dentro do buffer de destino. A principal vantagem sobre o laço de \texttt{MPI\_Recv} é a simplificação do código e a melhoria assintótica: a rotina coletiva implementa algoritmos de árvore internamente, reduzindo o custo de comunicação de $O(p)$ para $O(\log p)$.

% ---------------------------------------------------------
\subsection{4(b) Máximo e Mínimo Globais com \texttt{MPI\_Reduce}}

O processo 0 inicializou um vetor de $N = 10^7$ doubles e calculou os extremos serialmente como gabarito. A distribuição dos blocos foi feita com \texttt{MPI\_Scatter}. Cada processo calculou seu mínimo e máximo locais; em seguida, duas chamadas a \texttt{MPI\_Reduce} (com os operadores \texttt{MPI\_MIN} e \texttt{MPI\_MAX}) consolidaram os resultados globais no processo 0.

\textbf{Execução:}
\begin{lstlisting}[style=shell]
mpicc -O2 -Wall -o minmax minmax.c
mpiexec -n 2 ./minmax
mpiexec -n 4 ./minmax
\end{lstlisting}

\textbf{Saída:}
\begin{lstlisting}[style=shell]
--- Resultado Serial ---
Min = 0.0000000894 | Max = 0.9999999483
--- Resultado MPI ---
Min = 0.0000000894 | Max = 0.9999999483

--- Resultado Serial ---
Min = 0.0000000554 | Max = 0.9999999497
--- Resultado MPI ---
Min = 0.0000000554 | Max = 0.9999999497
\end{lstlisting}

Os resultados paralelos coincidem exatamente com o gabarito serial em ambos os casos. A utilização de \texttt{MPI\_Reduce} com \texttt{MPI\_MIN}/\texttt{MPI\_MAX} eliminou qualquer lógica explícita de comparação ou laços de \texttt{MPI\_Recv} no processo 0.

% =========================================================
\section{Exercício 5 — Soma Paralela de Vetores (Scatter / Gather)}

Foram implementadas duas versões. Em ambas, o processo 0 inicializa vetores \texttt{x[N]} (com $x_i = i$) e \texttt{y[N]} (com $y_i = 2i$), e \texttt{MPI\_Scatter} distribui blocos de \texttt{x} e \texttt{y} para todos os processos. Cada processo calcula seu bloco local $z = x + y$.

\begin{itemize}
  \item \textbf{Versão \texttt{mpi\_vecadd\_gather}:} \texttt{MPI\_Gather} coleta \texttt{z} apenas no processo 0, que verifica e imprime o resultado.
  \item \textbf{Versão \texttt{mpi\_vecadd\_allgather}:} \texttt{MPI\_Allgather} replica o \texttt{z} completo na memória de \emph{todos} os processos; cada rank confirma o recebimento e imprime.
\end{itemize}

\textbf{Execução:}
\begin{lstlisting}[style=shell]
make run
\end{lstlisting}

\textbf{Saída:}
\begin{lstlisting}[style=shell]
z = x + y:
  z[ 0] = 0 | z[ 1] = 3 | z[ 2] = 6 | ... | z[15] = 45

Rank 3 confirmando o recebimento completo. z = x + y:
  z[ 0] = 0 | z[ 1] = 3 | z[ 2] = 6 | ... | z[15] = 45
\end{lstlisting}

Ambas as versões produzem o vetor resultado correto ($z_i = 3i$).

\textbf{Quando \texttt{MPI\_Allgather} é preferível a \texttt{MPI\_Gather}?} A versão com \texttt{MPI\_Gather} é suficiente quando o único objetivo é consolidar o resultado no processo raiz para saída ou pós-processamento centralizado. Já \texttt{MPI\_Allgather} torna-se essencial em algoritmos onde o resultado global de uma etapa é o input necessário para \emph{todos} os processos na etapa seguinte. Um exemplo clássico é o produto matriz-vetor paralelo em métodos iterativos (e.g., gradiente conjugado), onde a matriz é particionada por blocos de linhas: para que cada processo multiplique suas linhas locais pela coluna inteira, ele precisa do vetor multiplicador completo. Utilizando \texttt{MPI\_Allgather}, os fragmentos locais são reconstruídos e replicados simultaneamente na memória de cada nó, sem a necessidade de um broadcast sequencial centralizado.

% =========================================================
\section{Exercício 6 — Tipos de Dados Derivados}

Foram implementadas duas versões para transmitir a estrutura:

\begin{lstlisting}[style=ccode]
struct Student { char name[50]; double grade; int id; };
\end{lstlisting}

\begin{itemize}
  \item \textbf{Versão \texttt{student\_struct}:} um tipo MPI derivado foi criado com \texttt{MPI\_Type\_create\_struct}, mapeando os três campos da struct. A macro \texttt{offsetof} de \texttt{<stddef.h>} foi usada para calcular os deslocamentos de memória de forma segura, garantindo que qualquer padding invisível inserido pelo compilador seja corretamente contabilizado. Após o registro com \texttt{MPI\_Type\_commit}, toda a estrutura foi transmitida com \textbf{uma única chamada} \texttt{MPI\_Bcast}.

  \item \textbf{Versão \texttt{student\_three\_bcasts}:} cada campo da estrutura foi transmitido separadamente com \texttt{MPI\_Bcast}, totalizando \textbf{três chamadas} de comunicação para um único registro.
\end{itemize}

\textbf{Execução:}
\begin{lstlisting}[style=shell]
make run
\end{lstlisting}

\textbf{Saída:}
\begin{lstlisting}[style=shell]
[Struct Bcast] Rank 0 recebeu -> Nome: Vivian Souza | Nota: 10.0 | ID: 123456789
[Struct Bcast] Rank 1 recebeu -> Nome: Vivian Souza | Nota: 10.0 | ID: 123456789
[Struct Bcast] Rank 2 recebeu -> Nome: Vivian Souza | Nota: 10.0 | ID: 123456789
[Struct Bcast] Rank 3 recebeu -> Nome: Vivian Souza | Nota: 10.0 | ID: 123456789

[3 Bcasts] Rank 0 recebeu -> Nome: Vivian Souza | Nota: 10.0 | ID: 123456789
[3 Bcasts] Rank 1 recebeu -> Nome: Vivian Souza | Nota: 10.0 | ID: 123456789
[3 Bcasts] Rank 2 recebeu -> Nome: Vivian Souza | Nota: 10.0 | ID: 123456789
[3 Bcasts] Rank 3 recebeu -> Nome: Vivian Souza | Nota: 10.0 | ID: 123456789
\end{lstlisting}

O conteúdo recebido é exatamente o mesmo em ambas as versões, confirmando que o mapeamento de memória via \texttt{MPI\_Type\_create\_struct} foi realizado corretamente. A diferença central está no número de chamadas de comunicação: \textbf{3 chamadas} na versão com três Bcasts contra \textbf{1 chamada} na versão com tipo derivado. Em ambientes distribuídos, cada invocação de uma rotina de comunicação paga uma latência fixa de rede (abertura de canal, criação de cabeçalhos, roteamento). A abordagem com três broadcasts paga essa latência três vezes por registro. Ao empacotar os campos heterogêneos em um único envelope via tipo derivado, a latência é paga apenas uma vez — um ganho que se multiplica significativamente quando esta rotina está inserida em laços iterativos de milhares de repetições, como em simulações numéricas e algoritmos de aprendizado de máquina distribuído.

\end{document}